\section{Scope and Review Method}\label{sec:methodology}\label{sec:background}
\subsection{Questions and eligibility}
We ask which hidden targets and observables appear (RQ1), how signal origin relates to receiver permission (RQ2), what establishes feasibility and transfer (RQ3), and which intervention addresses which observable at what cost (RQ4). An agent selects and executes actions using an LLM, tools, memory or orchestration; inherited serving and hardware channels need not be agent-specific.

Core eligibility requires (C1) LLM/agent execution, (C2) a hidden target, (C3) an execution-dependent auxiliary observable rather than ordinary output semantics alone, (C4) specified observation permissions, and (C5) an empirical inference evaluation. A controlled synthetic-state inference can satisfy C5; a proposed mechanism or timing gap alone cannot. Mixed papers qualify through an identifiable component. Table~\ref{tab:boundary-rules} supplies the decision rules.
\begin{table}[t]\centering\footnotesize
\caption{Observable-based eligibility. A hidden target alone does not make a semantic attack a side channel.}\label{tab:boundary-rules}
\begin{tabularx}{\columnwidth}{@{}p{.39\columnwidth}Y@{}}\toprule
Observed input & Decision and required distinction\\\midrule
Length, timing, cache response & Core when measured execution dependence supports hidden-target inference.\\
Own-site actions & Core attribution with site permission; no automatic access to other sites or private tool logs.\\
Reply states the secret & Direct disclosure; adjacent, not auxiliary-channel evidence.\\
Reply style/content identifies a user & Adjacent unless a separately evaluated auxiliary mechanism is established. WildAuth stays a mixed-view comparator.\\
Output-filter membership oracle & Can be core: the filter's execution, not ordinary semantic output inference, mediates the signal.\\
Privileged loop/fault logs & Separate diagnostic component when C1--C5 hold; does not increase confidentiality evidence.\\\bottomrule
\end{tabularx}\end{table}


The final coded collection contains \FinalN{} families available by October 8, 2026: \FinalCore{} core (\FinalA{} workflow-origin, \FinalB{} inherited-runtime/application) and \FinalAdjacent{} adjacent. Four core components concern authorized diagnostics. Other targets include content/property inference, provenance attribution and human/agent distinction; hiddenness alone establishes no confidentiality violation. Families are deduplicated studies; result records distinguish conditions within a study.

The auxiliary-only rule is applied consistently. Coding fingerprints and WildAuth remain adjacent because their evaluated views mix semantic and interaction features; feature importance is not a separately evaluated auxiliary-only attack \cite{ghaleb2026codingfingerprint,wang2026dialogue}. Metric-only fault localization qualifies as a diagnostic component despite operator authorization. ZeroPrint's response differences, endpoint-content fingerprints and StepFinder's mixed view remain adjacent \cite{scheinert2026rca,shao2026zeroprint,leshin2026endpoint,zhu2026stepfinder}.

\subsection{Executable scope and acquisition limits}
Native IEEE/ACM exports, a fixed 40-seed one-hop citation route and exact-title/DOI source lookup form the acquisition routes. Appendix~\ref{app:search} prints the executed queries, fields, dates and export rules. Source versions are merged only with identity evidence. Backward tracing covers 2,237 reference occurrences; forward lookup resolves 38 seeds and returns 24 distinct works. Neither route is recursively exhaustive.

We distinguish the acquired candidate frames from the coded collection. The prepared metadata frame contains 4,753 unique IDs with a recorded exclusion, referral or unresolved disposition. A bounded source reconciliation accounts for 109 overlapping decision rows in four nominated frames (57 follow-ups, six proposals, 22 backward candidates and 24 forward records). A disposition map includes pending material; it is not 109 independent full readings. Recent empirical channels, evaluated interventions, new receivers, prior systematizations and potential counterexamples receive source-review priority.

Broad database coverage remains incomplete: IEEE Q4's 35,743 historical hits were not fully acquired, and Q4b's 847 exports do not establish retained coverage. The dated ACM frame has 13,201 groups, including 883 referrals and 2,950 unresolved routes. The frames overlap and do not constitute a complete PRISMA funnel. Known-paper export yields are reported separately from indexing and field eligibility; they are not recall estimates. These limitations can conceal counterexamples or closer frameworks, not merely distort prevalence. Our conclusions therefore identify the exact result pairs that support them and expose single-source dependence.

\subsection{Coding, adjudication and reproducibility}
The extraction dictionary specifies receiver permissions, auxiliary versus semantic observations, result splitting, execution mode, utility/cost object and source/run binding. Source-scoped \NRR{}, unavailable \NRA{} and unverified \NRV{} remain distinct from replication. Code inspection can establish implementation behavior without authenticating the historical run.

Screening and targeted rechecks are AI-assisted. The 23-case label-blinded exercise has 17/21 agreement among determinate second judgments; the 332-dispute readbacks and separate exclusion audits are reported in Appendix~\ref{app:search}. They neither establish independent researcher coding nor estimate missed eligible studies. The supplied independent-review packets cover central sources, all mitigation families, unresolved defense assertions, boundary decisions and high-risk exclusions; no researcher returns have been received. Consequently, interpretation-sensitive judgments retain this validation limitation.

All \FinalCore{} core families have a scoped intervention disposition. The companion Supplementary Evidence package supplies the 60-family index, 33 result records, 66 utility/cost field positions, source versions/locators, claim links and a fit-free verifier. Table~\ref{tab:result-lineage} gives a complete chain in the paper itself. Appendix~\ref{app:claim-sensitivity} tests source deletion and unresolved bindings. The package can be reviewed as an accompanying file; public hosting and independent researcher validation are not claimed.
